\section{Exponential Fourier Series}

\subsubsection*{Motivation and Intuition}

The trigonometric Fourier series uses sine and cosine functions. However, using Euler's formula,
\[
    e^{i\theta}=\cos\theta+i\sin\theta,
\]
we can rewrite Fourier series in terms of complex exponentials.

This form is especially powerful in signal processing, communications, and systems analysis.

\begin{conceptbox}
    For a function $f(x)$ with period $L$, the exponential Fourier series is
    \[
        f(x)=\sum_{n=-\infty}^{\infty}
        c_n \exp\left(i\frac{2\pi n x}{L}\right),
    \]
    where
    \[
        c_n=\frac{1}{L}\int_0^L
        f(x)\exp\left(-i\frac{2\pi n x}{L}\right)\,dx,
        \qquad n\in\mathbb{Z}.
    \]
\end{conceptbox}

\subsubsection*{Interpretation}

The coefficient $c_n$ measures the contribution of the complex exponential with frequency
\[
    \nu_n=\frac{n}{L}.
\]

Positive values of $n$ correspond to positive frequency components, while negative values of $n$ correspond to negative frequency components.

For real-valued signals,
\[
    c_{-n}=\overline{c_n}.
\]

\subsubsection*{Relationship with Trigonometric Coefficients}

\begin{conceptbox}
    For $n\geq 1$, the trigonometric coefficients and exponential coefficients are related by
    \[
        c_n=\frac{a_n-ib_n}{2},
    \]
    \[
        c_{-n}=\frac{a_n+ib_n}{2},
    \]
    and
    \[
        c_0=\frac{a_0}{2}.
    \]
    Equivalently,
    \[
        a_n=c_n+c_{-n},
    \]
    and
    \[
        b_n=i(c_n-c_{-n}).
    \]
\end{conceptbox}

\subsubsection*{Worked Example}

\begin{examplebox}
    Find the exponential Fourier series coefficients of
    \[
        f(x)=\cos\left(\frac{2\pi x}{L}\right).
    \]

    Using Euler's formula,
    \[
        \cos\left(\frac{2\pi x}{L}\right)
        =
        \frac{1}{2}
        \left[
            \exp\left(i\frac{2\pi x}{L}\right)
            +
            \exp\left(-i\frac{2\pi x}{L}\right)
            \right].
    \]

    Compare this with
    \[
        f(x)=\sum_{n=-\infty}^{\infty}
        c_n\exp\left(i\frac{2\pi n x}{L}\right).
    \]

    The term
    \[
        \exp\left(i\frac{2\pi x}{L}\right)
    \]
    corresponds to $n=1$, so
    \[
        c_1=\frac{1}{2}.
    \]

    The term
    \[
        \exp\left(-i\frac{2\pi x}{L}\right)
    \]
    corresponds to $n=-1$, so
    \[
        c_{-1}=\frac{1}{2}.
    \]

    All other coefficients are zero:
    \[
        c_n=0,\qquad n\neq \pm 1.
    \]

    Therefore,
    \[
        \cos\left(\frac{2\pi x}{L}\right)
        =
        \frac{1}{2}\exp\left(i\frac{2\pi x}{L}\right)
        +
        \frac{1}{2}\exp\left(-i\frac{2\pi x}{L}\right).
    \]
\end{examplebox}

\subsubsection*{Engineering Insight}

\begin{noteBox}
    The exponential Fourier series is widely used in electrical engineering because linear systems act simply on complex exponentials. If an input has Fourier coefficients $c_n$, then a linear time-invariant system modifies each frequency component independently.
\end{noteBox}

\subsubsection*{Common Mistakes and Notes}

\begin{conceptbox}
    Common mistakes:
    \begin{enumerate}
        \item Forgetting that the exponential series sums from $n=-\infty$ to $\infty$.
        \item Confusing the period $L$ with $2L$. In the exponential form above, the period is $L$.
        \item Forgetting the negative sign in the coefficient formula:
              \[
                  c_n=\frac{1}{L}\int_0^L f(x)e^{-i2\pi nx/L}\,dx.
              \]
    \end{enumerate}
\end{conceptbox}